\section{Dépendances, procédures et reproduction}

Les constructions utilisées sont développées dans l’article :
le registre dodécaphasique dans la section~\ref{sec:k-registro}, la clôture
positionnelle dans la section~\ref{sec:alpha}, l’action dans la
section~\ref{hmtII:sec:action} et la paire angulaire dans la
section~\ref{hmtII:sec:ii-angulos}. La réponse constitutive de la
section~\ref{iii:sec:constitutiva} compose ces sorties en conservant
les étiquettes de feuillet et d’orientation.

\subsection{Ordre de reproduction}

La reproduction distingue le producteur de coordonnées des évaluateurs
ultérieurs. Le premier énumère les états admissibles et applique les règles
d’APP, de TRIT et de TPK. Ses sorties régionales sont des entrées légitimes des
opérateurs ultérieurs du même développement. Le programme d’évaluation
du vide n’est pas présenté comme un substitut de ce producteur.

L’ordre de lecture du paquet est : définitions du noyau ; énumération
et raffinement régional ; registre signé et transformation de $K$ ;
clôture positionnelle de $\alpha$ ; section d’action et retour ; paire angulaire ;
canaux contractifs ; opérateur constitutif ; moments cycliques et fonctionnelle
de Barbero ; sections dimensionnelles et comparaison ultérieure.

Chaque étape conserve son domaine. Un fichier de chiffres régionaux porte
l’empreinte du producteur et non une attribution métrologique. Le tuple de douze
composantes de $K$ a son propriétaire dans le chapitre du registre. Les
coefficients d’action, l’incidence $90/120$ et les multiplicités
$12:-1$ se lisent dans leurs démonstrations antérieures à l’évaluation. Le
programme local reproduit ces compositions avec une précision déclarée.

\subsection{Trois vérifications de fonctions différentes}

La conservation documentaire vérifie l’identité des fichiers et la continuité
des copies. La preuve mathématique établit les identités et les limites
dans leurs domaines. L’exécution numérique vérifie une évaluation concrète,
sa stabilité et les résidus des identités. Aucune de ces trois
vérifications ne remplace les autres.

Le paquet joint comprend un manifeste des sources, les textes LaTeX,
les données régionales utilisées et les programmes de vérification. Les preuves
d’inversion, de monotonie, de moments, de transport et de covariance se trouvent
dans le corps de l’article. Les sources documentaires sont citées par
provenance ; le lecteur peut vérifier les arguments développés ici
sans utiliser un résultat numérique cible comme générateur.

La clôture positionnelle de $\alpha$ est démontrée dans la
proposition~\ref{prop:alpha-normalizacion}. Le
théorème~\ref{thm:alpha-dos-publicaciones-completas} identifie ses
préfixes à ceux de la représentation analytique corrélée du même
état, au moyen de la récurrence explicite et des cylindres compatibles.
L’inversion constitutive de
la section~\ref{iii:sec:constitutiva} est formulée dans la normalisation
interne indiquée : les changements d’unités sont annulés avant d’appliquer
cette inverse. Ces conditions délimitent les résultats qu’utilisent
les compositions ultérieures.
